\documentclass[a4paper,12pt,twoside]{memoir}

% Castellano
\usepackage[spanish,es-tabla]{babel}
\selectlanguage{spanish}
\usepackage[utf8]{inputenc}
\usepackage[T1]{fontenc}
\usepackage{lmodern} % scalable font
\usepackage{microtype}
\usepackage{placeins}

\RequirePackage{booktabs}
\RequirePackage[table]{xcolor}
\RequirePackage{xtab}
\RequirePackage{multirow}

% Links
\PassOptionsToPackage{hyphens}{url}\usepackage[colorlinks]{hyperref}
\hypersetup{
	allcolors = {red}
}

% Ecuaciones
\usepackage{amsmath}

% Rutas de fichero / paquete
\newcommand{\ruta}[1]{{\sffamily #1}}

% Párrafos
\nonzeroparskip

% Huérfanas y viudas
\widowpenalty100000
\clubpenalty100000

% Evitar solapes en el header
\nouppercaseheads

% Imagenes
\usepackage{graphicx}
\newcommand{\imagen}[2]{
	\begin{figure}[!h]
		\centering
		\includegraphics[width=0.9\textwidth]{#1}
		\caption{#2}\label{fig:#1}
	\end{figure}
	\FloatBarrier
}

\newcommand{\imagenflotante}[2]{
	\begin{figure}%[!h]
		\centering
		\includegraphics[width=0.9\textwidth]{#1}
		\caption{#2}\label{fig:#1}
	\end{figure}
}



% El comando \figura nos permite insertar figuras comodamente, y utilizando
% siempre el mismo formato. Los parametros son:
% 1 -> Porcentaje del ancho de página que ocupará la figura (de 0 a 1)
% 2 --> Fichero de la imagen
% 3 --> Texto a pie de imagen
% 4 --> Etiqueta (label) para referencias
% 5 --> Opciones que queramos pasarle al \includegraphics
% 6 --> Opciones de posicionamiento a pasarle a \begin{figure}
\newcommand{\figuraConPosicion}[6]{%
  \setlength{\anchoFloat}{#1\textwidth}%
  \addtolength{\anchoFloat}{-4\fboxsep}%
  \setlength{\anchoFigura}{\anchoFloat}%
  \begin{figure}[#6]
    \begin{center}%
      \Ovalbox{%
        \begin{minipage}{\anchoFloat}%
          \begin{center}%
            \includegraphics[width=\anchoFigura,#5]{#2}%
            \caption{#3}%
            \label{#4}%
          \end{center}%
        \end{minipage}
      }%
    \end{center}%
  \end{figure}%
}

%
% Comando para incluir imágenes en formato apaisado (sin marco).
\newcommand{\figuraApaisadaSinMarco}[5]{%
  \begin{figure}%
    \begin{center}%
    \includegraphics[angle=90,height=#1\textheight,#5]{#2}%
    \caption{#3}%
    \label{#4}%
    \end{center}%
  \end{figure}%
}
% Para las tablas
\newcommand{\otoprule}{\midrule [\heavyrulewidth]}
%
% Nuevo comando para tablas pequeñas (menos de una página).
\newcommand{\tablaSmall}[5]{%
 \begin{table}
  \begin{center}
   \rowcolors {2}{gray!35}{}
   \begin{tabular}{#2}
    \toprule
    #4
    \otoprule
    #5
    \bottomrule
   \end{tabular}
   \caption{#1}
   \label{tabla:#3}
  \end{center}
 \end{table}
}

%
%Para el float H de tablaSmallSinColores
\usepackage{float}

%
% Nuevo comando para tablas pequeñas (menos de una página).
\newcommand{\tablaSmallSinColores}[5]{%
 \begin{table}[H]
  \begin{center}
   \begin{tabular}{#2}
    \toprule
    #4
    \otoprule
    #5
    \bottomrule
   \end{tabular}
   \caption{#1}
   \label{tabla:#3}
  \end{center}
 \end{table}
}

\newcommand{\tablaApaisadaSmall}[5]{%
\begin{landscape}
  \begin{table}
   \begin{center}
    \rowcolors {2}{gray!35}{}
    \begin{tabular}{#2}
     \toprule
     #4
     \otoprule
     #5
     \bottomrule
    \end{tabular}
    \caption{#1}
    \label{tabla:#3}
   \end{center}
  \end{table}
\end{landscape}
}

%
% Nuevo comando para tablas grandes con cabecera y filas alternas coloreadas en gris.
\newcommand{\tabla}[6]{%
  \begin{center}
    \tablefirsthead{
      \toprule
      #5
      \otoprule
    }
    \tablehead{
      \multicolumn{#3}{l}{\small\sl continúa desde la página anterior}\\
      \toprule
      #5
      \otoprule
    }
    \tabletail{
      \hline
      \multicolumn{#3}{r}{\small\sl continúa en la página siguiente}\\
    }
    \tablelasttail{
      \hline
    }
    \bottomcaption{#1}
    \rowcolors {2}{gray!35}{}
    \begin{xtabular}{#2}
      #6
      \bottomrule
    \end{xtabular}
    \label{tabla:#4}
  \end{center}
}

%
% Nuevo comando para tablas grandes con cabecera.
\newcommand{\tablaSinColores}[6]{%
  \begin{center}
    \tablefirsthead{
      \toprule
      #5
      \otoprule
    }
    \tablehead{
      \multicolumn{#3}{l}{\small\sl continúa desde la página anterior}\\
      \toprule
      #5
      \otoprule
    }
    \tabletail{
      \hline
      \multicolumn{#3}{r}{\small\sl continúa en la página siguiente}\\
    }
    \tablelasttail{
      \hline
    }
    \bottomcaption{#1}
    \begin{xtabular}{#2}
      #6
      \bottomrule
    \end{xtabular}
    \label{tabla:#4}
  \end{center}
}

%
% Nuevo comando para tablas grandes sin cabecera.
\newcommand{\tablaSinCabecera}[5]{%
  \begin{center}
    \tablefirsthead{
      \toprule
    }
    \tablehead{
      \multicolumn{#3}{l}{\small\sl continúa desde la página anterior}\\
      \hline
    }
    \tabletail{
      \hline
      \multicolumn{#3}{r}{\small\sl continúa en la página siguiente}\\
    }
    \tablelasttail{
      \hline
    }
    \bottomcaption{#1}
  \begin{xtabular}{#2}
    #5
   \bottomrule
  \end{xtabular}
  \label{tabla:#4}
  \end{center}
}



\definecolor{cgoLight}{HTML}{EEEEEE}
\definecolor{cgoExtralight}{HTML}{FFFFFF}

%
% Nuevo comando para tablas grandes sin cabecera.
\newcommand{\tablaSinCabeceraConBandas}[5]{%
  \begin{center}
    \tablefirsthead{
      \toprule
    }
    \tablehead{
      \multicolumn{#3}{l}{\small\sl continúa desde la página anterior}\\
      \hline
    }
    \tabletail{
      \hline
      \multicolumn{#3}{r}{\small\sl continúa en la página siguiente}\\
    }
    \tablelasttail{
      \hline
    }
    \bottomcaption{#1}
    \rowcolors[]{1}{cgoExtralight}{cgoLight}

  \begin{xtabular}{#2}
    #5
   \bottomrule
  \end{xtabular}
  \label{tabla:#4}
  \end{center}
}




\graphicspath{ {./img/} }

% Capítulos
\chapterstyle{bianchi}
\newcommand{\capitulo}[2]{
	\setcounter{chapter}{#1}
	\setcounter{section}{0}
	\setcounter{figure}{0}
	\setcounter{table}{0}
	\chapter*{#2}
	\addcontentsline{toc}{chapter}{#2}
	\markboth{#2}{#2}
}

% Apéndices
\renewcommand{\appendixname}{Apéndice}
\renewcommand*\cftappendixname{\appendixname}

\newcommand{\apendice}[1]{
	%\renewcommand{\thechapter}{A}
	\chapter{#1}
}

\renewcommand*\cftappendixname{\appendixname\ }

% Formato de portada
\makeatletter
\usepackage{xcolor}
\newcommand{\tutor}[1]{\def\@tutor{#1}}
\newcommand{\course}[1]{\def\@course{#1}}
\definecolor{cpardoBox}{HTML}{E6E6FF}
\def\maketitle{
  \null
  \thispagestyle{empty}
  % Cabecera ----------------
\noindent\includegraphics[width=\textwidth]{cabecera}\vspace{1cm}%
  \vfill
  % Título proyecto y escudo informática ----------------
  \colorbox{cpardoBox}{%
    \begin{minipage}{.8\textwidth}
      \vspace{.5cm}\Large
      \begin{center}
      \textbf{TFG del Grado en Ingeniería Informática}\vspace{.6cm}\\
      \textbf{\LARGE\@title{}}
      \end{center}
      \vspace{.2cm}
    \end{minipage}

  }%
  \hfill\begin{minipage}{.20\textwidth}
    \includegraphics[width=\textwidth]{escudoInfor}
  \end{minipage}
  \vfill
  % Datos de alumno, curso y tutores ------------------

%  \begin{center}%
%  {%
%    \noindent\LARGE
%    Presentado por \@author{}\\ 
%    en Universidad de Burgos --- \@date{}\\
%    Tutores: \@tutor{}\\
%  }%
%  \end{center}%
\begin{center}%
   {%
    \noindent\LARGE
    Presentado por \@author{}\\ 
    en Universidad de Burgos \\
    el \@date{}\\
    \begin{tabbing}
    Tutores: \= \kill
    Tutores: César García Osorio\\
           \> Alicia Olivares Gil\\
    \end{tabbing}
  }%
  \end{center}%
  \null
  \cleardoublepage
  }
\makeatother


% Datos de portada
\title{AI-Music}
\author{José Ángel López Estrada}
\tutor{César Ignacio García Osorio, Alicia Olivares Gil}
\date{\today}

\begin{document}

\maketitle



\cleardoublepage



%%%%%%%%%%%%%%%%%%%%%%%%%%%%%%%%%%%%%%%%%%%%%%%%%%%%%%%%%%%%%%%%%%%%%%%%%%%%%%%%%%%%%%%%



\frontmatter


\clearpage

% Indices
\tableofcontents

\clearpage

\listoffigures

\clearpage

\listoftables

\clearpage

\mainmatter

\appendix

\apendice{Plan de Proyecto Software}

\section{Introducción}

Este plan de proyecto de software se refiere al desarrollo de una aplicación que clasifica automáticamente los géneros musicales basándose en características de audio. 
Este proyecto ha sido concebido como un trabajo de fin de grado y ha implicado la implementación de técnicas de aprendizaje automático y procesamiento de señales de audio digitales para lograr el objetivo.
La aplicación busca proporcionar una solución efectiva para clasificar pistas musicales de forma automática.

En este apéndice se describirá la planificación sobre la que se ha desarrollado el proyecto.

\section{Planificación temporal}

El desarrollo del proyecto se organizó siguiendo la metodología ágil de SCRUM, con iteraciones semanales o "<sprints">. 
Cada sprint implicó una serie de actividades que culminaron con una entrega incremental del proyecto.

\textbf{Sprints totales}: 14

\subsection{Sprint 1 (08/03/2023 - 15/03/2023)}

En el Sprint 1, se realizó la configuración inicial del proyecto, añadiendo una descripción detallada en el archivo \texttt{README.md} y estableciendo las bases de la documentación. Asimismo, se implementaron componentes básicos y se trabajó en el diseño de la aplicación, aprovechando las funcionalidades de la biblioteca \textit{librosa} para el procesamiento de audio.

\subsection{Sprint 2 (15/03/2023 - 22/03/2023)}

En el Sprint 2, se configuró el entorno de desarrollo utilizando Poetry, que posteriormente se descartó ante el uso de entornos virtuales clásicos, y se implementó el módulo \texttt{train\textunderscore classifier} para entrenar el clasificador de géneros musicales. De igual manera, también se llevó a cabo la extracción de características MFCC y se realizó un preprocesamiento básico de los datos.

\subsection{Sprint 3 (22/03/2023 - 05/04/2023)}

En el Sprint 3, se realizaron pequeños cambios en el código e investigación sobre el área.

\subsection{Sprint 4 (05/04/2023 - 12/04/2023)}

En el Sprint 4, se realizaron diversas tareas. En primer lugar, se agregó un archivo \texttt{requirements.txt} para gestionar las dependencias del proyecto y se creó la estructura inicial de Flask para desarrollar la aplicación web. Complementariamente, se diseñó una aplicación Flask básica y se llevó a cabo el preprocesamiento de géneros musicales. Además, se realizaron ajustes en la estructura del proyecto y se actualizó la extracción de características MFCC. También se implementó la visualización de características de audio, como la forma de onda, el espectrograma, el cromagrama y los MFCC. Finalmente, se realizaron mejoras en el estilo y se añadió un modelo de aprendizaje automático inicial.

\subsection{Sprint 5 (12/04/2023 - 19/04/2023)}

En el Sprint 5, se trabajó en los aspectos visuales del proyecto. Se crearon mockups e imágenes de prueba para la interfaz de usuario, permitiendo una mejor comprensión de cómo se presentaría la aplicación.

\subsection{Sprint 6 (19/04/2023 - 26/04/2023)}

En el Sprint 6, se avanzó en la sección de carga del proyecto y se añadió una página de inicio. Además, se implementó un reproductor de música que permitía a los usuarios escuchar las pistas seleccionadas.

\subsection{Sprint 7 (26/04/2023 - 10/05/2023)}

En el Sprint 7, se llevaron a cabo varias tareas importantes. Se añadió un nuevo modelo al sistema y se actualizó la documentación correspondiente. Paralelamente, también se trabajó en la implementación de las páginas de registro e inicio de sesión para los usuarios, aunque finalmente no pudieron ser llevadas a cabo. También se realizaron ajustes en el guardado del modelo y se añadieron las páginas de inicio de sesión y registro al índice del proyecto.

\subsection{Sprint 8 (10/05/2023 - 17/05/2023)}

En el Sprint 8, se dedicó tiempo a la búsqueda de referencias y documentación relevante para el proyecto. Se recopilaron trabajos relacionados, papers e información general que proporcionaron una base sólida para la consolidación de la aplicación.

\subsection{Sprint 9 (17/05/2023 - 24/05/2023)}

En el Sprint 9, se realizaron mejoras en la extracción de características y se llevó a cabo un proceso de limpieza, eliminando archivos de borrador que ya no eran necesarios.

\subsection{Sprint 10 (24/05/2023 - 02/06/2023)}

En el Sprint 10, se añadió documentación adicional para complementar el proyecto, asegurando que todos los aspectos importantes estuvieran debidamente explicados.

\subsection{Sprint 11 (02/06/2023 - 09/06/2023)}

En el Sprint 11, se implementaron generadores para la fase de entrenamiento del modelo, lo que permitió una mejor gestión de los datos. También se actualizó la representación de géneros predichos.

\subsection{Sprint 12 (09/06/2023 - 16/06/2023)}

En el Sprint 12, se realizaron cambios en el diseño de la aplicación, mejorando la representación visual de las características de audio. Se llevaron a cabo actualizaciones en el diseño de los colores utilizados en la interfaz.

\subsection{Sprint 13 (16/06/2023 - 23/06/2023)}

En el Sprint 13, se realizaron varios cambios significativos. Por un lado se implementaron diferentes mejoras en el procesamiento y entrenamiento de los datos. De igual modo, se actualizaron el reproductor de audio y la información de las pistas. Además, se realizaron múltiples actualizaciones en la memoria, los anexos y la aplicación web. Eventualmente, se llevaron a cabo correcciones menores de estilo para mejorar la presentación general.

\subsection{Sprint 14 (23/06/2023 - 05/06/2023)}

En el Sprint 14, se continuaron realizando actualizaciones y mejoras en la documentación. Se realizaron cambios específicos en la sección B de los anexos, relacionada con los requisitos funcionales y no funcionales. También se actualizaron las secciones A, D y E de los anexos, que incluyen el plan del proyecto, el manual del programador y el manual de usuario, respectivamente. Se agregó un diagrama general de casos de uso y se llevaron a cabo actualizaciones adicionales en la documentación, la aplicación web y los anexos. Además se desplegó la aplicación web en un servidor en la nube, concretamente en la plataforma \texttt{PythonAnywhere}.

\section{Estudio de viabilidad}

\subsection{Viabilidad económica}

La viabilidad económica de un proyecto de software se refiere a la capacidad para generar ingresos tanto para cubrir el desarrollo como para proporcionar una rentabilidad a medio y largo plazo. El presente proyecto no está planteado con un objetivo económico, por lo que esta sección se va a dedicar a explorar un posible plan de monetización de forma ficticia, explorando los costes totales de desarrollo de software y potenciales métodos de monetización.

\subsection{Costes de hardware}
Los costes de hardware se refieren al gasto económico que se realiza para obtener los diferentes elementos hardware que se necesitan para poder llevar a cabo el desarrollo del proyecto. Los elementos hardware pueden ser ordenadores, periféricos o componentes internos por ejemplo.
En este caso se ha planteado la compra de un ordenador portátil de unos 1000€. Esta decisión se toma teníendo en cuenta los siguientes factores:

\begin{itemize}
\tightlist
\item \textbf{Rendimiento}: Este proyecto requiere un uso de algoritmos de aprendizaje automático y procesamiento de datos, por lo que se demanda una cierta capacidad de computación. 
Un equipo dentro del rango de precio planteado viene equipado con un procesador potente, como un Intel Core i7 (12th o 13th Gen) o un AMD Ryzen 7 (6th Gen). Además suelen incluir una tarjeta gráfica dedicada, como una NVIDIA RTX 3060 con 6 GB de VRam, con la ventaja en rendimiento que esto conlleva en el uso de bibliotecas de \textit{Deep Learning} como TensorFlow.

\item \textbf{Durabilidad}: Los ordenadores de este rango de precio suelen contar con componentes de mayor calidad asegurando una mayor vida útil. La durabilidad es importante en proyectos como este ya que un fallo de hardware podría tener un impacto significativo en la productividad.

\item \textbf{Soporte}: Los fabricantes proporcionan un soporte de mayor calidad en ordenadores situados en este rango de precio. Este soporte puede incluir mejores condiciones de garantía o actualizaciones más duraderas en partes vitales del sistema como la BIOS.
\end{itemize}

\begin{table}[h]
\centering
\begin{tabular}{|l|c|}
\hline
\textbf{Concepto} & \textbf{Coste} \\ 
\hline
Ordenador portátil & 1000€ \\ 
\hline
Total & 1000€ \\ 
\hline
\end{tabular}
\caption{Costes de hardware}
\end{table}

\subsection{Costes de software}
Los costes de software se refieren al gasto económico que se realiza en la adquisición de las distintas licencias de software que son necesarias para la realización del proyecto.

\begin{itemize}
\tightlist
\item \textbf{Sistema Operativo}: Se ha elegido el sistema operativo Windows 11 Home como base para realizar el desarrollo del proyecto. Se podrían plantear alternativas gratuitas como alguna distribución de Linux pero por sencillez y extensión de uso se ha elegido Windows.

\item \textbf{Servicios en la nube}: Durante las etapas intensivas de entrenamiento de modelos y almacenamiento de datos, podría llegar a ser necesario el uso de servicios de computación en la nube como AWS \cite{AWS} o Microsoft Azure \cite{Kcpitt}. El coste puede variar según el uso necesario. Por ejemplo en un proyecto de estas características, contando con unos 100 GB de datos en AWS, se estimaría un costo de almacenamiento de unos 2€/mensuales y un costo de procesamiento de unos 30€/mensuales, utilizando 10 horas en una instancia \texttt{p3.2xlarge}, la cual es comúnmente utilizada para labores de aprendizaje automático. En total, el coste de la computación en la nube a lo largo del proyecto ha sido de 192€.
\end{itemize}

\begin{table}[h]
\centering
\begin{tabular}{|l|c|}
\hline
\textbf{Concepto} & \textbf{Coste} \\
\hline
Windows 11 Home \cite{win11} & 145€ \\
Computación en la nube & 192€ \\
\hline
\textbf{Total} & \textbf{337€} \\
\hline
\end{tabular}
\caption{Costes de software para el desarrollo del proyecto}
\end{table}

\subsection{Otros costes}

\begin{table}[h]
\centering
\begin{tabular}{|l|c|}
\hline
\textbf{Concepto} & \textbf{Coste} \\
\hline
Electricidad (mensual) & 50€ \\
Internet (mensual) & 35€ \\
Espacio de trabajo (mensual) & 350€ \\
\hline
\textbf{Total (mensual)} & \textbf{435€} \\
\hline
\textbf{Total (6 meses)} & \textbf{2610€} \\
\hline
\end{tabular}
\caption{Otros costes para el desarrollo del proyecto}
\end{table}

\subsection{Costes totales}
Teniendo en cuenta una duración de desarrollo de proyecto de 6 meses, el coste total ha sido:

\begin{table}[h]
\centering
\begin{tabular}{|l|c|}
\hline
\textbf{Concepto} & \textbf{Coste} \\
\hline
Costes de hardware & 50€ \\
Costes de software & 32€ \\
Otros costes & 2610€ \\
\hline
\textbf{Total} & \textbf{2692€} \\
\hline
\end{tabular}
\caption{Costes totales de desarrollo del proyecto}
\end{table}

\subsection{Ingresos}
La idea del proyecto es generar ingresos hasta el punto de rentabilizar el gasto de desarrollo como mínimo. 
Existen diversas formas en las que la aplicación podría generar ingresos como por ejemplo:

\begin{itemize}
\tightlist

\item \textbf{Sistema de suscripción}: Se puede ofrecer la aplicación con distintas características según el \textit{tier} o el tipo de suscripción. Podrían existir tres \textit{tier}:
	\begin{itemize}
	\tightlist
	\item \textbf{Suscripción gratuita}: límite 5 canciones/diarias

	\item \textbf{Suscripción silver 5€/mes}: canciones ilimitadas.

	\item \textbf{Suscripción gold 15€/mes}: canciones ilimitadas y análisis con más detalle.
	\end{itemize}

\item \textbf{Publicidad}: Se puede monetizar el servicio a través de anuncios. Los usuarios que usen la versión gratuita de la aplicación verán anuncios en distintas partes de la interfaz, mientras que si pagan una mensualidad (5€/mes) tiene la opción de eliminar anuncios.

\item \textbf{Integraciones de API}: Se puede ofrecer un servicio de uso de la API a aplicaciones de terceros para utilizar la aplicación. Esta API puede ser gratuita, con limitaciones en el número de canciones a predecir, o de pago, con un menor número de limitaciones.
\end{itemize}

\subsection{Viabilidad legal}

La viabilidad legal del proyecto se refiere al cumplimiento de diversas leyes y obligaciones a la hora de desarrollar el software.

\subsection{Derechos de autor}

El tema de los derechos de autor en el mundo de la música es un tema complicado. Los derechos de autor de la música son un área legal compleja que envuelve partes como artistas, productores y discografías. Este proyecto va a utilizar un conjunto de datos musicales libres de derechos de autor lo que facilita el trabajo eliminando esta parte del proceso de desarrollo.

A la hora de la subida de archivos musicales para su clasificación existe un problema, el usuario final puede subir cualquier fichero de audio, tenga derechos de autor o no. Por lo tanto, el procedimiento a realizar será la eliminación de cualquier fichero musical de la aplicación tras realizar el proceso de predicción.

\subsection{Licencias de software}

Según el software utilizado, se tienen en cuenta las siguientes licencias \cite{Creative_Commons, Open_Source_Initiative_2023}:

\begin{table}[h]
\centering
\caption{Licencias de las herramientas utilizadas}
\begin{tabular}{|l|l|}
\hline
\textbf{Herramienta} & \textbf{Licencia} \\
\hline
Python \cite{3.10_Documentation} & PSF \\
Librosa \cite{librosa} & ISC \\
TensorFlow \cite{TensorFlow} & Apache 2.0 \\
ScikitLearn \cite{scikit} & BSD-3-Clause \\
Flask \cite{Flask} & BSD-3-Clause \\
FMA Dataset \cite{fma_dataset} & CC BY-NC-SA 4.0 \\
Matplotlib \cite{matplotlib} & PSF \\
Pandas \cite{pandas} & BSD-3-Clause \\
Numpy \cite{NumPy} & BSD-3-Clause \\
\hline
\end{tabular}
\end{table}

Según las herramientas utilizadas, se ha elegido la licencia MIT para la distribución y uso del software desarrollado. Esta es una licencia de código abierto permisiva, la cual permite a los usuarios utilizar, modificar y distribuir el software bajo ciertas condiciones, requiriendo la inclusión del aviso de derechos de autor y la renuncia de responsabilidad en el software distribuido.
Una de las características principales de la licencia MIT es su flexibilidad. Permite a los usuarios utilizar el software con pocas restricciones y sin imponer condiciones adicionales. Además, ofrece libertad para modificar y distribuir el software tanto en proyectos de código abierto como en proyectos comerciales.
La elección de la licencia MIT se basa en su adecuación a las necesidades y especificidades de este proyecto. Algunas razones por las cuales es la licencia que mejor se ajusta son:

\begin{enumerate}
\tightlist

\item \textbf{Flexibilidad y colaboración}: La licencia MIT \cite{Open_Source_Initiative_2023} permite que otras personas utilicen, modifiquen y distribuyan el software sin imponer restricciones excesivas. Esto fomenta la colaboración y la participación de la comunidad, lo cual es beneficioso para un proyecto de desarrollo de software.
\item \textbf{Compatibilidad}: La licencia MIT es compatible con la mayoría de las otras licencias de código abierto. Esto es importante, ya que el proyecto utiliza diversas herramientas con diferentes licencias, asegurando que no haya conflictos y permitiendo una integración fluida entre ellas.
\item \textbf{Reconocimiento y protección de derechos de autor}: La licencia MIT incluye disposiciones que protegen los derechos de autor y exigen la inclusión del aviso de derechos de autor en el software distribuido. Esto garantiza el reconocimiento adecuado del autor original del proyecto y protege sus derechos.

\end{enumerate}
\apendice{Especificación de Requisitos}

\section{Introducción}

Esta sección describe la especificación de requisitos para el proyecto. Se proporciona información detallada sobre los requisitos funcionales (RF) y requisitos no funcionales (RNF) así como el detalle de los casos de uso (CU).

\section{Objetivos generales}

\begin{itemize}
\tightlist
\item Desarrollar un sistema de reconocimiento de estilos musicales utilizando inteligencia artificial.
\item Diseñar e implementar una aplicación web que permita usar el modelo de una forma sencilla.
\item Obtener conclusiones y conocimiento a partir de los datos.
\end{itemize}

\section{Catalogo de requisitos}

\subsection{Requisitos funcionales}

Un requisito funcional es una especificación que describe lo que un sistema debe hacer o cómo debe comportarse.
Los requisitos funcionales pueden incluir detalles como cálculos, manipulación de datos, interacción con el usuario, etc.
A continuación se listan los requisitos funcionales extraidos para este proyecto:

\begin{itemize}
\tightlist

\item \textbf{RF-1 Predicción de estilo musical}: la aplicación debe ser capaz de predecir el estilo musical de una canción.
	\begin{itemize}
	\tightlist
	\item \textbf{RF-1.1 Extracción de características de audio}: la aplicación debe ser capaz de extraer características relevantes de la canción que se utilizarán para predecir el estilo musical.

	\item \textbf{RF-1.2 Implementar un modelo de machine learning para realizar la predicción}: la aplicación debe implementar un modelo de machine learning que se entrenará con las características de un conjunto de datos musical.
	\end{itemize}

\item \textbf{RF-2 Presentación de la información}: la aplicación debe mostrar información relevante sobre la predicción realizada.
	\begin{itemize}
	\tightlist
	\item \textbf{RF-2.1 Detalles de la predicción}: la aplicación debe ser capaz de mostrar datos sobre los detalles de la predicción como la confianza (probabilidad) o posibles géneros musicales alternativos.

	\item \textbf{RF-2.2 Información de la canción}: la aplicación debe mostrar información relevante sobre la pista de audio o canción utilizada como la duración, formato o frecuencia de muestreo.
	\end{itemize}
\item \textbf{RF-3 Implementación de una interfaz de usuario}: la aplicación debe contener una interfaz de usuario para facilitar su uso.
	\begin{itemize}
	\tightlist
	\item \textbf{RF-3.1 Funcionalidad de carga de archivos}: la interfaz de usuario debe incluir una función de carga de archivos que permita seleccionar y cargar una canción para realizar la predicción.

	\item \textbf{RF-3.2 Representación de los resultados}: la interfaz de usuario debe incluir sección donde se presenten de forma clara los resultados de la predicción.
	\end{itemize}

\end{itemize}
\subsection{Requisitos no funcionales}

Los requisitos no funcionales se centran en las características del sistema que no están directamente relacionadas con su comportamiento, por ejemplo rendimiento o seguridad.
A continuación se listan los requisistos no funcionales extraídos para este proyecto:

\begin{itemize}
\tightlist
\item \textbf{RNF-1 Rendimiento}: la aplicación debe ser capaz de procesar y analizar una canción y devolver los resultados dentro de un tiempo razonable, idealmente unos pocos segundos.

\item \textbf{RNF-2 Seguridad}: la aplicación debe cumplir con protocolos de seguridad como cifrado de comunicaciones o la eliminación de la canción después de su procesamiento.

\item \textbf{RNF-3 Compatibilidad}: la aplicación debe ser compatible con diversos navegadores y sistemas operativos.

\item \textbf{RNF-4 Usabilidad}: la aplicación debe ser fácil de usar para el usuario. El usuario debería ser capaz de realizar el proceso de subida de archivos y comprobar la predicción y resultados de una forma sencilla.

\item \textbf{RNF-5 Portabilidad}: la aplicación debe ser compatible con diversos dispositivos como PCs, tablets o dispositivos móviles.
\end{itemize}

\section{Especificación de requisitos}

Los casos de uso son descripciones detalladas del funcionamiento de una parte del sistema desde la perspectiva del usuario.

En este sección se mostrará el diagrama general de casos de uso y se explicará en detalle cada uno de ellos.

\subsection{Diagrama general de casos de uso}

\imagen{Diagrama_general_de_casos_de_uso}{Diagrama general de casos de uso.}

\subsection{Casos de uso}

A continuación se exponen las especificaciones de los casos de uso que forman la aplicación.

% Caso de Uso 1 -> Carga de la canción
\begin{table}[p]
	\centering
	\begin{tabularx}{\linewidth}{ p{0.21\columnwidth} p{0.71\columnwidth} }
		\toprule
		\textbf{CU-1}    & \textbf{Carga de la canción}\\
		\toprule
		\textbf{Versión}              & 1.0    \\
		\textbf{Autor}                & José Ángel López \\
		\textbf{Requisitos asociados} & RF-1.1, RF-3.1 \\
		\textbf{Descripción}          & El usuario utiliza la función de carga de archivo para seleccionar y cargar una canción en la aplicación. \\
		\textbf{Precondiciones}        & 
		\begin{enumerate}		    
			\def\labelenumi{\arabic{enumi}.}
			\tightlist
			\item La aplicación se encuentra en funcionamiento.
			\item El usuario tiene acceso a los archivos de música que desea cargar.
		\end{enumerate}\\
		\textbf{Acciones}             &
		\begin{enumerate}
			\def\labelenumi{\arabic{enumi}.}
			\tightlist
			\item El usuario selecciona la opción para cargar una canción pulsando el botón \textit{Upload File}.
			\item El usuario busca y selecciona la canción que desea cargar desde su almacenamiento local.
			\item La canción se carga en la aplicación.
		\end{enumerate}\\
		\textbf{Postcondiciones} &
		\begin{enumerate}
			\def\labelenumi{\arabic{enumi}.}
			\tightlist
			\item La canción se encuentra cargada en la aplicación.
			\item La aplicación está lista para extraer características y realizar el proceso de predicción de la canción cargada.
		\end{enumerate}\\
		\textbf{Excepciones} &
		\begin{enumerate}
			\def\labelenumi{\arabic{enumi}.}
			\tightlist
			\item El archivo cargado no es un fichero de audio.
			\item La canción tiene un formato incompatible.
		\end{enumerate}\\
		\textbf{Importancia}          & Alta \\
		\bottomrule
	\end{tabularx}
	\caption{CU-1 Carga de la canción}
\end{table}

% Caso de Uso 2 -> Extracción de características de audio.
\begin{table}[p]
	\centering
	\begin{tabularx}{\linewidth}{ p{0.21\columnwidth} p{0.71\columnwidth} }
		\toprule
		\textbf{CU-2}    & \textbf{Extracción de características de audio.}\\
		\toprule
		\textbf{Versión}              & 1.0    \\
		\textbf{Autor}                & José Ángel López \\
		\textbf{Requisitos asociados} & RF-1.1 \\
		\textbf{Descripción}          & La aplicación extrae las características relevantes del audio para su posterior análisis. \\
		\textbf{Precondiciones}        & 
		\begin{enumerate}		    
			\def\labelenumi{\arabic{enumi}.}
			\tightlist
			\item La canción se ha cargado correctamente en la aplicación.
		\end{enumerate}\\
		\textbf{Acciones}             &
		\begin{enumerate}
			\def\labelenumi{\arabic{enumi}.}
			\tightlist
			\item La aplicación inicia el proceso de extracción de características del archivo de audio cargado.
		\end{enumerate}\\
		\textbf{Postcondiciones} &
		\begin{enumerate}
			\def\labelenumi{\arabic{enumi}.}
			\tightlist
			\item Las características relevantes de la canción se han extraído correctamente.
			\item La aplicación está lista para utilizar las características extraídas para realizar la predicción del estilo musical.
		\end{enumerate}\\
		\textbf{Excepciones} &
		\begin{enumerate}
			\def\labelenumi{\arabic{enumi}.}
			\tightlist
			\item Ocurrió un error durante la extracción de las características del audio.
		\end{enumerate}\\
		\textbf{Importancia}          & Alta \\
		\bottomrule
	\end{tabularx}
	\caption{CU-2 Extracción de características de audio}
\end{table}

% Caso de Uso 3 -> Ejecución del modelo de machine learning
\begin{table}[p]
	\centering
	\begin{tabularx}{\linewidth}{ p{0.21\columnwidth} p{0.71\columnwidth} }
		\toprule
		\textbf{CU-3}    & \textbf{Ejecución del modelo de machine learning}\\
		\toprule
		\textbf{Versión}              & 1.0    \\
		\textbf{Autor}                & José Ángel López \\
		\textbf{Requisitos asociados} & RF-1.2 \\
		\textbf{Descripción}          & La aplicación ejecuta el modelo de machine learning para realizar la predicción del estilo musical según las caracterísiticas extraídas. \\
		\textbf{Precondiciones}        & 
		\begin{enumerate}		    
			\def\labelenumi{\arabic{enumi}.}
			\tightlist
			\item Las características de audio de la canción se han extraído correctamente.
			\item Las características de audio de la canción son suficientes para realizar el proceso de predicción.
		\end{enumerate}\\
		\textbf{Acciones}             &
		\begin{enumerate}
			\def\labelenumi{\arabic{enumi}.}
			\tightlist
			\item El modelo procesa las características extraídas de la canción.
			\item Se obtiene el resultado de la predicción, el cual consiste en un vector de probabilidades entre diferentes estilos musicales.
		\end{enumerate}\\
		\textbf{Postcondiciones} &
		\begin{enumerate}
			\def\labelenumi{\arabic{enumi}.}
			\tightlist
			\item La aplicación ha generado una predicción sobre el estilo musical de la canción.
			\item La aplicación está lista para mostrar los detalles de la predicción al usuario.
		\end{enumerate}\\
		\textbf{Excepciones} &
		\begin{enumerate}
			\def\labelenumi{\arabic{enumi}.}
			\tightlist
			\item El modelo no puede procesar las características extraídas.
		\end{enumerate}\\
		\textbf{Importancia}          & Alta \\
		\bottomrule
	\end{tabularx}
	\caption{CU-3 Ejecución del modelo de machine learning}
\end{table}

% Caso de Uso 4 -> Visualización de la predicción
\begin{table}[p]
	\centering
	\begin{tabularx}{\linewidth}{ p{0.21\columnwidth} p{0.71\columnwidth} }
		\toprule
		\textbf{CU-4}    & \textbf{Visualización de la predicción}\\
		\toprule
		\textbf{Versión}              & 1.0    \\
		\textbf{Autor}                & José Ángel López \\
		\textbf{Requisitos asociados} & RF-2.1, RF-2.2 \\
		\textbf{Descripción}          & Tras la ejecución del modelo de machine learning, la aplicación muestra la predicción del estilo musical al usuario. \\
		\textbf{Precondiciones}        & 
		\begin{enumerate}		    
			\def\labelenumi{\arabic{enumi}.}
			\tightlist
			\item La aplicación ha generado una predicción.
		\end{enumerate}\\
		\textbf{Acciones}             &
		\begin{enumerate}
			\def\labelenumi{\arabic{enumi}.}
			\tightlist
			\item La aplicación muestra la predicción del estilo musical en la interfaz.
			\item El usuario puede cambiar entre distintas visualizaciones de forma dinámica.
			\item El usuario visualiza la predicción.
		\end{enumerate}\\
		\textbf{Postcondiciones} &
		\begin{enumerate}
			\def\labelenumi{\arabic{enumi}.}
			\tightlist
			\item El usuario ha obtenido la predicción del estilo musical de la canción.
		\end{enumerate}\\
		\textbf{Excepciones} &
		\begin{enumerate}
			\def\labelenumi{\arabic{enumi}.}
			\tightlist
			\item Se produce un error al intentar mostrar la predicción.
		\end{enumerate}\\
		\textbf{Importancia}          & Alta \\
		\bottomrule
	\end{tabularx}
	\caption{CU-4 Visualización de la predicción}
\end{table}

% Caso de Uso 5 -> Reproducción de la canción
\begin{table}[p]
	\centering
	\begin{tabularx}{\linewidth}{ p{0.21\columnwidth} p{0.71\columnwidth} }
		\toprule
		\textbf{CU-5}    & \textbf{Reproducción de la canción}\\
		\toprule
		\textbf{Versión}              & 1.0    \\
		\textbf{Autor}                & José Ángel López \\
		\textbf{Requisitos asociados} & RF-1.2, RF-1.3 \\
		\textbf{Descripción}          & La aplicación permite al usuario reproducir la canción cargada. \\
		\textbf{Precondiciones}        & 
		\begin{enumerate}		    
			\def\labelenumi{\arabic{enumi}.}
			\tightlist
			\item La canción ha sido cargada correctamente en la aplicación.
		\end{enumerate}\\
		\textbf{Acciones}             &
		\begin{enumerate}
			\def\labelenumi{\arabic{enumi}.}
			\tightlist
			\item El usuario selecciona la opción de reproducir la canción mediante el botón \textit{play} del reproductor.
			\item La aplicación reproduce la canción cargada.
		\end{enumerate}\\
		\textbf{Postcondiciones} &
		\begin{enumerate}
			\def\labelenumi{\arabic{enumi}.}
			\tightlist
			\item La canción está siendo reproducida en la aplicación.
		\end{enumerate}\\
		\textbf{Excepciones} &
		\begin{enumerate}
			\def\labelenumi{\arabic{enumi}.}
			\tightlist
			\item Se produce un error en la reproducción de la canción.
		\end{enumerate}\\
		\textbf{Importancia}          & Media \\
		\bottomrule
	\end{tabularx}
	\caption{CU-5 Reproducción de la canción}
\end{table}

% Caso de Uso 6 -> Carga de una nueva canción
\begin{table}[p]
	\centering
	\begin{tabularx}{\linewidth}{ p{0.21\columnwidth} p{0.71\columnwidth} }
		\toprule
		\textbf{CU-6}    & \textbf{Carga de una nueva canción}\\
		\toprule
		\textbf{Versión}              & 1.0    \\
		\textbf{Autor}                & José Ángel López \\
		\textbf{Requisitos asociados} & RF-1.1, RF-1.3 \\
		\textbf{Descripción}          & El usuario puede cargar una nueva canción sin tener que cerrar o reiniciar la aplicación. \\
		\textbf{Precondiciones}        & 
		\begin{enumerate}		    
			\def\labelenumi{\arabic{enumi}.}
			\tightlist
			\item La aplicación ha realizado la predicción de una canción.
		\end{enumerate}\\
		\textbf{Acciones}             &
		\begin{enumerate}
			\def\labelenumi{\arabic{enumi}.}
			\tightlist
			\item El usuario hace selecciona el nombre de la aplicación, en la zona superior izquierda de la pantalla, para volver a la página inicial.
			\item El usuario pulsa el botón \textit{Upload File}.
			\item El usuario busca y selecciona la nueva canción que desea cargar desde su almacenamiento.
		\end{enumerate}\\
		\textbf{Postcondiciones} &
		\begin{enumerate}
			\def\labelenumi{\arabic{enumi}.}
			\tightlist
			\item La nueva canción se carga en la aplicación.
		\end{enumerate}\\
		\textbf{Excepciones} &
		\begin{enumerate}
			\def\labelenumi{\arabic{enumi}.}
			\tightlist
			\item El archivo cargado no es un fichero de audio.
			\item La canción tiene un formato incompatible.
		\end{enumerate}\\
		\textbf{Importancia}          & Alta \\
		\bottomrule
	\end{tabularx}
	\caption{CU-6 Carga de una nueva canción}
\end{table}
\apendice{Especificación de diseño}

\section{Introducción}

Cómo se han resuelto e implementado los requisitos anteriormente mencionados. Definición de datos, estructuras, directorios, clases, etc.

\section{Diseño de datos}

Tipos de datos con los que se van a trabajar.

INSERTAR DIAGRAMA DE CLASES, RELACIONES, E/R, ETC.

\section{Diseño procedimental}

DISEÑO PROCEDIMENTAL SIGNIFICA CÓMO SE HAN DISEÑADO LOS PROCEDIMIENTOS DE LA APLICACIÓN (CÓMO LOS MÉTODOS, FUNCIONES, ETC.)
MOSTRAR CADA PROCEDIMIENTO EN PSEUDOCODIGO O LENGUAJE NATURAL.
diagramas de interacción, secuencia, etc.

\section{Diseño arquitectónico}

MODELO ARQUITECTONICO PREVISTO, PATRONES DE DISEÑO. Por ejemplo decir si hemos usado patrón vista controlador, patrón MVM, etc.
herencias, genericidad, clases, etc.

EXPLICAR LA FORMA GENERAL DE LA ARQUITECTURA DE LA APLICACIÓN (CON DIAGRAMAS, FLECHAS, ETC.)

INSERTAR NUEVA SECCIÓN CON DISEÑO DE INTERFACES SI PROCEDE.
\apendice{Documentación técnica de programación}

\section{Introducción}
Este apéndice aborda la documentación técnica de programación de la aplicación. En las siguientes secciones se detallan diferentes detalles importantes para trabajar en el proyecto, como pueden ser:
\begin{enumerate}
\tightlist

\item Estructura de directorios.
\item Instalación del entorno de desarrollo.
\item Obtención del código fuente.
\item Ejecución de la aplicación.
\item Pruebas realizadas.

\end{enumerate}

\section{Estructura de directorios}

La estructura de directorios de la aplicación es la siguiente:

\begin{itemize}
\tightlist

\item \texttt{/}: Directorio raíz del proyecto. Contiene el archivo \texttt{requirements.txt} con la información de las dependencias del proyecto, el archivo \texttt{README.md} con información relevante sobre el proyecto y el archivo \texttt{LICENSE} con la licencia del proyecto.

\item \texttt{/app/}: Directorio con el código fuente de la aplicación. Además incluye en fichero \texttt{\textunderscore\textunderscore init\textunderscore\textunderscore .py} que lanza y define las rutas en la aplicación web utilizando \texttt{Flask}.

\item \texttt{/app/models/}: Directorio donde se almacenan los modelos de machine learning entrenados que la aplicación utiliza para realizar las predicciones de los estilos musicales.

\item \texttt{/app/src/}: Directorio donde se encuentra el código fuente Python. Aquí es donde se realiza todo el procesamiento interno de la aplicación.

\item \texttt{/app/static/}: Directorio donde se almacenan los archivos estáticos utilizados en la interfaz web de la aplicación. Por ejemplo, archivos CSS, pistas de audio subidas o scripts de JavaScript.

\item \texttt{/app/templates/}: Directorio donde se almacenan los archivos HTML que la aplicación utiliza para generar las páginas web.

\item \texttt{/app/tests/}: Directorio donde se almacenan los tests unitarios y de integración utilizados para verificar el correcto funcionamiento de la aplicación.

\item \texttt{/data/}: Directorio donde se almacenan los datos utilizados por la aplicación para realizar el entrenamiento o el procesado.

\item \texttt{/data/processed/}: Directorio donde se almacenan los datos procesados y listos para ser entrenados en el modelo de machine learning.

\item \texttt{/data/raw/}: Directorio donde se almacenan los ficheros de audio sin procesar.

\item \texttt{/docs/latex/}: Documentación del proyecto en formato \LaTeX.

\item \texttt{/examples/}: Ejemplos de uso de las librerías.

\end{itemize}

\section{Manual del programador}

Esta sección tiene como objetivo ser una guía para que los futuros programadores puedan entender el código fuente, configurar el entorno de desarrollo y poder contribuir en el proyecto.

Para el desarrollo del proyecto se utilizó Visual Studio Code con el lenguaje de programación Python. Además, se utilizó Git para el sistema de control de versiones. Se recomienda tener ciertos conocimientos en estos sistemas antes de empezar.

\subsection{Herramientas recomendadas}

Para contribuir al desarrollo del proyecto se recomienda utilizar las siguientes herramientas.

\begin{itemize}
\tightlist

\item \textbf{Git}: Se trata del sistema de control de versiones utilizado para realizar un seguimiento de los cambios en el código fuente de la aplicación a lo largo del tiempo. Se puede combinar con GitHub. \cite{Git} \cite{GitHub}

\imagen{Git_console}{Consola de Git.}

\item \textbf{Visual Studio Code}: Es el editor de código fuente recomendado para continuar con el desarrollo del proyecto. Se trata de un editor de código simple y ligero pero con una gran cantidad de extensiones que facilitan el desarrollo.
Por ejemplo, incluye linting de código, formateo automático del código y soporte con Jupyter Notebooks o Docker. \cite{Microsoft_2021}

\imagen{VSCode}{Visual Studio Code.}

\item \textbf{Jupyter Notebooks}: Herramienta de desarrollo de código y visualizaciones. Es una herramienta muy útil para realizar experimentación ya que une la ejecución de código fuente con la visualización de datos y la documentación.
No se trata de una herramienta para realizar una versión final del producto, pero es muy útil para probar y generar prototipos de una forma rápida. \cite{Jupyter}

\imagen{Jupyter_notebook}{Jupyter Notebook.}

\end{itemize}

\subsection{Explicación de aspectos relevantes de programación}

\section{Compilación, instalación y ejecución del proyecto}

Para obtener el código fuente del proyecto los pasos a seguir son los siguientes:

\begin{itemize}
\tightlist

\item \textbf{Clonar el repositorio Git}
	\begin{itemize}
	\tightlist
		\item \texttt{git clone https://github.com/jle1001/AI-Music.git}
	\end{itemize}

\item \textbf{Crear un entorno virtual de Python}: 
	\begin{itemize}
	\tightlist
		\item \texttt{python3 -m venv venv}
	\end{itemize}

\item \textbf{Activar el entorno virtual}: 
	\begin{itemize}
		\item Linux: \texttt{source venv/bin/activate}
		
		\item Windows: \texttt{venv\textbackslash Scripts\textbackslash activate}
	\end{itemize}

\item \textbf{Instalación de dependencias}
	\begin{itemize}
	\tightlist
		\item \texttt{pip install -r requirements.txt}
	\end{itemize}

\item \textbf{Iniciar el servidor}
	\begin{itemize}
	\tightlist
		\item \texttt{flask run --debug}
	\end{itemize}
\end{itemize}

Una vez que el servidor está en funcionamiento, se puede acceder a la aplicación web a través de un navegador entrando a la dirección \texttt{localhost:5000}.

\section{Pruebas del sistema}

\subsection{Pruebas unitarias}

Las pruebas unitarias son un tipo de pruebas de software en las que se comprueban individualmente las partes más pequeñas del sistema. Estas partes pueden ser funciones, métodos o clases.

\subsection{Pruebas de integración}

Las pruebas de integración son una fase de pruebas de software en la que se prueban varias partes del programa para verficiar que interactúan correctamente entre ellas. Estas pruebas se centran en la interfaz y en la interacción entre las diferentes partes.

\subsection{Pruebas de interfaz de usuario}

Las pruebas de interfaz de usuario son pruebas que se utilizan para verificar que la interacción entre el usuario y el software se realiza correctamente.
\apendice{Documentación de usuario}

\section{Introducción}

Esta sección proporciona instrucciones de uso detalladas sobre cómo utilizar la aplicación, desde el proceso de instalación hasta el proceso de subida de archivos.

\section{Requisitos de usuarios}

\subsection{Entorno local}

\begin{table}[h]
	\centering
	\begin{tabular}{|c|c|}
		\hline
		\textbf{Requisito} & \textbf{Especificación}\\
		\hline
		CPU & 2.0 GHz o superior\\
		\hline
		RAM & 4GB o superior\\
		\hline
		Disco duro & 2GB de almacenamiento libre\\
		\hline
		Sistema Operativo & Compatible con Python 3.10 o superior\\
		\hline
		Python & Versión 3.10 o superior\\
		\hline
		Bibliotecas de Python & Especificadas en la sección de instalación\\
		\hline
	\end{tabular}
	\caption{Requisitos para el Funcionamiento Local del Sistema}
\end{table}

\begin{itemize}
\tightlist

\item \textbf{CPU 2.0 GHz o superior}: Un procesador de 2 GHz es capaz de manejar las instrucciones por segundo necesarias para realizar las tareas requeridas en esta aplicación. Como el procesamiento de audio y la ejecución de la predicción.

\item \textbf{Memoria RAM de 4GB o superior}: Los 4GB de RAM son necesarios para que la aplicación funcione de manera eficiente sin ralentizaciones o bloqueos.

\item \textbf{2GB de almacenamiento libre}: Esta cantidad de almacenamiento libre es la mínima necesaria para instalar la aplicación y sus librerías necesarias.

\item \textbf{Sistema operativo compatible con Python 3.10 o superior}: La aplicación está escrita en Python, por lo que se requiere un sistema operativo que pueda soportar la versión de Python 3.10 o superior.

\end{itemize}

\section{Instalación y ejecución}

Para utilizar la aplicación se puede acceder al siguiente enlace: \url{https://jle1001.pythonanywhere.com}

\section{Manual del usuario}
El objetivo de este manual es proporcionar una guía paso a paso para que el usuario pueda utilizar las funciones de la aplicación.

\subsection{Subida de canciones}
\begin{itemize}
\tightlist

\item En la ventana principal de la aplicación hay que pulsar el botón \textit{Upload File}.
\item Elegir el fichero de audio en el almacenamiento.
\item Una vez subido el fichero de audio se darán opciones como: reproducir la canción, visualizar datos básicos y elegir el modelo para la predicción.

\end{itemize}

\imagen{upload_file}{Subida de canciones.}

\subsection{Análisis y predicción}

\begin{itemize}
\tightlist

\item En la ventana principal de la aplicación, pulsar el botón \textit{Upload File}.
\item Elegir el fichero de audio en el almacenamiento.
\item Una vez subido el fichero de audio se darán opciones como: reproducir la canción, visualizar datos básicos y elegir el modelo para la predicción.
\item Para realizar la predicción del género musical se debe elegir un modelo a utilizar y pulsar el botón \textit{Predict}.

\end{itemize}

\imagen{info_cancion}{Información sobre la canción subida.}

\imagen{seleccion_algoritmo}{Selección de algoritmo.}

\subsection{Visualización del resultado}

\begin{itemize}
\tightlist

\item En la ventana principal de la aplicación, pulsar el botón \textit{Upload File}.
\item Elegir el fichero de audio en el almacenamiento.
\item Una vez subido el fichero de audio se darán opciones como: reproducir la canción, visualizar datos básicos y elegir el modelo para la predicción.
\item Para realizar la predicción del género musical se debe elegir un modelo a utilizar y pulsar el botón \textit{Predict}.
\item Una vez el modelo haya terminado de predecir el estilo musical, este se mostrará en la parte superior de la pantalla.
\item En la parte central de la pantalla se podrán ver cuatro gráficos con diferentes características musicales.

\end{itemize}

\imagen{waveform}{Visualización del resultado: Onda de audio.}

\imagen{spectrogram}{Visualización del resultado: Espectrograma.}

\imagen{chromagram}{Visualización del resultado: Cromograma}

\imagen{MFCC}{Visualización del resultado: MFCC.}

\subsection{Videos explicativos}

Demostración de uso de la aplicación: \href{https://www.youtube.com/watch?v=aX596dmo6Dk}{AI-Music - Demostración de funcionamiento}


\bibliographystyle{IEEEtran}
\bibliography{bibliografiaAnexos}

\end{document}
